\begin{abstract}
For cubic graphs $G$ and $H$, an $H$-coloring of $G$ is a map from the edges of $G$ to the edges
of $H$ that sends every vertex star of $G$ bijectively onto a vertex star of $H$; a Petersen
coloring is the case where $H$ is the Petersen graph. Ma, Mattiolo, Steffen and Wolf showed that
there is a unique minimal set $\mathcal{H}_3$ of connected bridgeless cubic graphs that color
every bridgeless cubic graph, and that a graph belongs to it exactly when no bridgeless cubic
graph of smaller order colors it. Since the Petersen coloring conjecture was refuted in 2026,
$\mathcal{H}_3$ is infinite and contains every smallest counterexample. Goedgebeur, Jooken,
M\'a\v{c}ajov\'a, Mattiolo, Mazzuoccolo and Ulyanov asked whether their two counterexamples on 52
vertices, the smallest known, are colorable only by themselves. We prove two lemmas on the vertex
map of an $H$-coloring: all its fibers have the same parity, and target vertices that are not
used can be reduced to at most one by the splitting lemma. They leave three kinds of colorings,
in none of which the target graph has a part unconstrained by $G$, and they turn the question
into one propositional formula per target order, with the target graph as part of the unknown.
For both 52-vertex counterexamples, every even target order from 40 to 50 is refuted with a DRAT
proof checked by \texttt{drat-trim}; several smaller orders are refuted as well. Together with the
result of Goedgebeur et al.\ that every bridgeless cubic graph on at most 38 vertices has a
Petersen coloring, this shows that both graphs belong to $\mathcal{H}_3$: they are colorable
only by themselves. Without the two lemmas the same search decided no target order of the first
graph within comparable time.
\end{abstract}
